\chapter{بحث، نتیجه‌گیری و پیشنهادها}

\section{بحث در مورد یافته‌های کلیدی}

این بخش با ترکیب نتایج اصلی و تشریح اهمیت آن‌ها، یافته‌های مستخرج از رساله را مورد ارزیابی عمیق قرار می‌دهد.

\subsection{نظریه‌پردازی مکانیزم: چرا \lr{RS-DRL} کار می‌کند؟}

همان‌طور که در نتایج ارزیابی تجربی (مندرج در فصل ۵ و با استناد به شکل‌های ۵ و ۶ در مقاله مرجع \cite{RS_DRL_Article}) نشان داده شد، یکی از مهم‌ترین نقاط قوت \lr{RS-DRL} توانایی آن در فرار از بهینه‌های محلی است که یک چالش جدی و متداول در سیستم‌های \lr{DRL} محسوب می‌شود. در محیط‌های پویا، پاداش متوسط مدل \lr{RS-DRL} نه تنها به‌طور پیوسته افزایش می‌یابد، بلکه در وضعیت پایداری نسبت به رویکرد پایه \lr{$\epsilon$-greedy DQN} - که با نوسانات شدیدی همراه است و نمی‌تواند از راه‌حل‌های غیربهینه رهایی یابد - در سطحی بسیار بالاتر تثبیت می‌گردد.

اثربخشی رویکرد \lr{RS-DRL} را می‌توان به سه مکانیزم به هم‌پیوسته زیر نسبت داد:

\begin{enumerate}
  \item \textbf{اکتشاف از طریق تزریق پاداش خوش‌بینانه:} با تغییر تدریجی و تصادفیِ تجربیات ناموفق گذشته ($r_i=0$) به تجربیات موفق ($r_i=1$)، عامل یادگیرنده سیگنال‌های تقویت‌کننده مثبتی را از اعمال و وضعیت‌هایی دریافت می‌کند که در شرایط عادی از اولویت خارج می‌شدند. این امر از همگرایی زودهنگام مدل به سیاست‌های غیربهینه که یکی از معضلاتِ بنیادین \lr{DQN} استاندارد است، ممانعت به عمل می‌آورد.
  \item \textbf{شکستن همبستگیِ زمانی در الگوهای شکست:} انتخاب تصادفی در فرآیند تغییر تجربیات مردود (که توسط پارامتر $\rho$ کنترل می‌گردد) متغیری را معرفی می‌کند که مانع از حفظ بیش‌حد (\lr{Overfitting}) عامل بر روی شرایط متوالی شکست می‌شود. این متغیر به ویژه در محیط‌های غیرایستایی (\lr{Non-stationary}) که در آن‌ها شکست دیروز ممکن است استراتژی بهینه امروز تلقی گردد، بسیار راهگشا خواهد بود.
  \item \textbf{یادگیری مستقل از دامنه (\lr{Domain-Agnostic}):} برخلاف روش‌های متکی به \lr{HER} که مستلزم تعاریف صریح و قطعی از اهداف هستند، شکل‌دهی تصادفیِ \lr{RS-DRL} بدون هیچ‌گونه نیاز به درک معناییِ اهدافِ دامنه عمل می‌کند. این موضوع، اجرایِ آن را ذاتاً برای سیستم‌هایی انعظاف‌پذیر می‌سازد که با اهداف تلویحی و آستانه‌محور - که در سیستم‌های خودوفقی بسیار شایع می‌باشند - سروکار دارند.
\end{enumerate}

\subsection{تحلیل تعاملات و مصالحه‌ها (به‌عنوان مثال، تاخیر در مقابل از دست رفتن بسته)}

در جداول مقایسه‌ای پیشین (برای نمونه جدول ۶ در توصیفات ارزیابی)، مشاهده گردید که برای تخصیص هدف \lr{TTS} در شبکه \lr{DeltaIoTv1}، مدل \lr{RS-DRL} به درستی به تاخیر متوازنی دست پیدا کرد (مقدار ۰.۰۹۷ در برابر ۰.۱۰۰ روش مرجع) و در مجموع عملکرد بسیار بهتری در مقایسه با \lr{DLASER+} نشان داد. با این وجود، برای شبکه \lr{DeltaIoTv2} تحت هدف \lr{TT} مقدار تاخیر بالاتری حاصل شد (۰.۲۵۱ در برابر ۰.۱۰۰ روش مرجع).

افزایش تاخیر در نسخه دوم شبکه برای مقاصد مربوط به \lr{TT} نشان‌دهنده یک مصالحه (\lr{Trade-off}) بسیار مهم در عملکرد عامل است: هنگامی‌که هدف کلیدیِ تعریف‌شده از سوی کاربر (مثلاً حفظ توانایی توان عملیاتی شبکه) با اولویت‌بندی به حداقل رساندن مشکل فقدانِ بسته توأم گردد، مدلِ هوشمند درمی‌یابد که باید به منظور جلوگیری از دست‌افتادگیِ پیام‌ها، به سمتِ مسیرهای قابل اعتمادتری متمایل شود که به صورت ذاتی کندتر هستند. در محیط‌های بهینه‌سازی چندهدفه، چنین مصالحه‌ای نه تنها یک نقطه ضعف و محدودیت به شمار نمی‌رود، بلکه از قضا یک رفتار برآمده و مطلوب از سوی عامل ارزیابی می‌گردد. چرا که در اینجا اولویت سیستم بر تضمینِ رساندن بسته‌ها بیش از حفظ سرعت بوده است و در کاربردهایی که صحت اطلاعات در آن‌ها اهمیت حیاتی دارد (مانند مانیتورینگ زیست‌محیطی یا سیستم‌های کنترل صنعتی) این تصمیم به‌عنوان عاقلانه‌ترین انتخاب مهندسی تلقی خواهد شد.

این رفتار کاملاً با اصول انتخاب بهینه در استراتژی‌های پرتو-اپتیمال (\lr{Pareto Optimality}) همخوانی دارد. هنگامی‌که مرزهای متضاد انتخاب‌ها بررسی می‌گردند، تنظیمات دارای تأخیرِ بی‌نهایت کم‌تر عموماً ارتباط معناداری با نرخ بالای اتلاف بسته‌ها – به دلیل استراتژی‌های تهاجمی‌ترِ انتقال داده – نشان می‌دهند. راهکار \lr{RS-DRL} یاد می‌گیرد این خطِ مصالحه را با ظرافت طی کند و تصمیم‌ها را منطبق بَر بستر نیاز و با تخصیص وزن‌های اختصاص‌یافته برای هر متغیر اخذ نماید.

\subsection{پیامدهای مهندسی سیستم‌های خودوفقی}

ثبات، پایداری و همگرایی سریع متد \lr{RS-DRL} عوامل حیاتی در محیط‌های پویایی هستند که مقیاس و شرایط آن‌ها به سرعت در حال جابجایی است. از طریق در هم تنیدگی مکانیزم شکل‌دهی پاداش، \lr{RS-DRL} تضمین می‌نماید که عامل تصمیم‌ساز حتی بدون نیاز به بازآموزی و تخصیص پیکربندی‌های جدید به تغییرات پاسخ می‌دهد، در حالی که در اکثر \lr{DRL}های استاندارد الزام به طی مجدد مسیر آموزش برای اداره کردن امور تغییر یافته، مشهود و گریز‌ناپذیر است. این خصوصیت، این مدل را برای سامانه‌های برخط مانند سیستم‌های ابری و شبکه‌های اشیاء که استقرار و بروزرسانی خودکار استراتژی در آن‌ها واجب است، بسیار متمایز می‌سازد.

به بیان دقیق‌تر، پیامدهای عملی این رهیافت برای مهندسین توسعه‌دهنده محیط‌های خودوفقی شامل موارد زیر است:

\begin{enumerate}
  \item \textbf{کاهش قابل‌ملاحظه سربار عملیاتی:} مدل‌های سنتی به هنگام برخورد با رانش و تغییر شرایط محیطی نیازمند صرف چرخه‌های متعدد بازآموزیِ پیاپی می‌باشند. در حالی که ویژگی‌های تعمیم‌پذیری بالای \lr{RS-DRL} امکان دور ماندنِ مستمرِ مدل از بروزرسانی‌های مجدد را فراهم می‌سازد که خود علاوه بر صرفه‌جویی شدید محاسباتی باعث جلوگیری از عملکرد ضعیف مدل در طی ادوارِ گذارِ بازآموزی می‌گردد.
  \item \textbf{بی‌نیازی از داده‌های برچسب‌دار:} بر خلاف مدل‌هایی همانند \lr{ML4EAS} یا \lr{DLASER+} که اتکای گسترده‌ای به تجمیعِ پیشینِ داده‌ها داشته‌اند، الگوریتم \lr{RS-DRL} کاملاً آنلاین و درلحظه توسعه می‌یابد. این قابلیت نقطه قوتِ برجسته‌ای در محیط‌هایی محسوب می‌گردد که تهیه اطلاعات پیش‌آموزشی غیرممکن و پیش‌بینی‌ناپذیر است.
  \item \textbf{پایداری و کاهش تخریب تدریجیِ کیفیت متناسب با گسترش:} استواریِ معنادارِ کیفیت و عملکرد در مواجهه با نسخه‌ی ۲۰۱۶-گزینه‌ای (نسخه \lr{v1})، گسترش ۱۲۹۶-گزینه‌ای (نسخه \lr{v1.1}) و مقیاس سرسام‌آور ۴۰۹۶ حالت ممکن (نسخه \lr{v2}) مؤید آن است که این راهکار متناسب با میزان گستردگی شبکه قابل استقرار است. همچنین افزایش مقطع زمانی همگرایی مدل، به نسبتِ پیچیدگیِ فضای جستجو، شیبی کمتر از حالت خطی (\lr{Sublinear}) دارد.
  \item \textbf{تصمیم‌گیریِ شفاف و قابل مدل‌سازی:} رویه‌هایِ رفتارگونه‌ی عامل (همانند پدیده مصالحه در برخورد با اتلافِ دیتا و زمانِ توقف) به نحو کاملاً معنا‌دار و قابل تعریفی مدل‌گستر و قابل توضیح می‌باشد که این ویژگی به مهندسان قدرت تحلیل مضاعف و قابلیت اطمینانِ بیشتری عرضه می‌نماید.
\end{enumerate}

\section{محدودیت‌ها و تهدیدات اعتبار (\lr{Threats to Validity})}

مانند تمامی‌رهیافت‌های پژوهشی، متدولوژی معرفی‌شده با محدودیت‌هایی مواجه است که بررسی صریحِ روایی و صحتِ نتایجِ استخراج‌شده، بستر مناسب را برای پژوهش‌های مکملِ آینده تبیین می‌نماید.

\subsection{اعتبار داخلی (\lr{Internal Validity})}

\begin{enumerate}
  \item \textbf{حساسیت نسبت به ابرپارامترها:} فاکتورِ کنترل‌کننده‌ی شکل‌دهیِ تجربیات تصادفی ($\rho$)، متولی تنظیم نمودن درصدِ اختصاص تجربه‌هایِ ناموفقِ تغییرِ یافته به نمونه‌های فرضیِ مثبت و امیدوارکننده است. با اینکه از طریق کاوش شبکه‌ای جامع (\lr{Grid Search}) توانستیم مناسب‌ترین مقادیر ممکن برای بهینه‌سازیِ عملکرد یافت کنیم، لیکن این مقادیر می‌توانند به نسبت سایر سیستم‌ها و دامنه‌های متفاوت، دستخوش نوسان باشند. بدیهی‌ست چنانچه پارامتر بسیار بالا وضع گردد موجب اختلال در روند یادگیری گردیده و مقادیر ضعیفِ آن نیز سبب نقصان در عامل فشارِ اکتشاف خواهد گشت.
  \item \textbf{نوسانات بذر تصادفیِ تولید ارقام (\lr{Random Seed Sensitivity}):} با وجود اجرای بیش از ۳۰ دور فرآیندِ آموزش مستقل همراه با صحت‌سنجی مستحکم آماری، ماهیتِ درگیر با احتمالات - چه در پویایی خود محیط و چه در بدنه تنظیم پاداش‌ها - سبب می‌شود تا خروجی‌های هر دَوره آموزشی با درجاتِ کوچکی از انحراف واریانس مواجه گردند. با این وصف، فواصل و حاشیه‌های اطمینان به دست آمده (قید شده در جداول فصل ۴) تبیین‌گر کمیت این نوسانات هستند.
  \item \textbf{مفروضات بازده دودویی (\lr{Binary Reward Assumption}):} یکی از شرایط پیش‌نیازِ به کار رفته در الگوریتم، اتکا بر مفروضات مرتبط با نتایج قطعی پیامدهایِ انطباقی به‌عنوان اهدافِ توفیق/سقوطِ مطلقِ دودویی (یا در مقیاس‌های استاندارد مشابه) است که نمی‌توان آن را به سهولت با کلیه نیازهای سیستم‌های ترکیبیِ غول‌پیکر وفق داد.
\end{enumerate}

\subsection{اعتبار خارجی (\lr{External Validity})}

\begin{enumerate}
  \item \textbf{تعمیم‌پذیری ساختار و تخصیصِ حوزه (\lr{Domain Specificity}):} فرآیند تمرکز و سنجشِ نهاییِ ارزیابی این رساله منحصراً روی بستر محیط شبکه‌های متراکم و سنسورهای اینترنتی (\lr{DeltaIoT}) انجام شد. هرچند پلتفرم مذکور آینه‌ای از مشکلات اساسی سیستم‌های مدرن را به تصویر کشیده و تاییدِ کیفیت کرده است، با این وصف، اعمال الگوریتمِ توسعه‌یافته بر دیگر شبکه‌هایِ توزیع‌یافته، نیازمند اعتبارسنجیِ مستقلی خواهد بود.
  \item \textbf{تمرکز بهینه بر شبکه‌های گسسته:} تدوین و صورت‌بندی فعلی به نحوی طراحی گشته که متناسب با فضای تصمیم‌داریِ مدل‌های محدود گسسته پایه‌ریزی شود. لیکن اکثریت سیستم‌های پویا (مانند سیستم‌های مقیاس‌بندی و تنظیمِ ظرفیتِ پردازشی‌ ابری) از پارامترهای پیوسته استفاده می‌نمایند و تنظیم معماری \lr{RS-DRL} متناسب با متدهای کلاسیک پیوسته نظیر شبکه‌های \lr{DDPG} ساختارهای بسط یافته مطالعاتی جدیدی را می‌طلبد.
  \item \textbf{توسعه‌ی در حوزه‌هایِ حساسِ ایمنی:} استفاده از استراتژی خوش‌بینی بر روی اقداماتی که ذاتاً و در دنیای فیزیکی خطرساز بوده‌اند (همانند پلتفرم‌های هدایت اتومبیل‌های هوشمند و یا سیستم‌های دارورسانی)، بدون در نظر گرفتن شرط احتیاط، ممکن است باعث ترغیب ماشین به انجام فعالیت‌های منجر به حادثه گردد و افزودن حلقه‌های تکمیلی تایید صلاحیت بسیار با اهمیت است.
\end{enumerate}

\subsection{اعتبار سازه (\lr{Construct Validity})}

\begin{enumerate}
  \item \textbf{شفافیت متریک‌ها:} سنجش عملکرد عامل هوش از قبیل بررسی مجانبیِ ثبات کیفیت زمانی تا رسیدن به راندمان مطلوب آستانه، رویه‌هایی پذیرفته شده می‌باشند؛ اما معرفی فاکتورهای سنجشی جایگزین مانند نرخ رنجوری (رگرت) و شاخص پایداری ممکن است ابعاد ناشناخته دیگری را نمایان سازند.
  \item \textbf{اصطکاک سنجشی مدل‌های جایگزین:} از آنجا که مدل‌های برگرفته از \lr{HER} نیازمند تنظیم صریحِ برچسب‌گذاریِ مقصد بوده‌اند، مقایسه صورت پذیرفته نیاز به در نظر گرفتن جایگزین‌های ضمنی کیفی داشت. هرچند مستندسازی اعمال شده بسیار متوازن بود ولی باید اذعان داشت که استفاده از این مقادیرِ تعدیل یافته معادل بهینه‌ای برای تمام ظرفیت موجود این دست از رقبا نمی‌تواند قلمداد گردد.
\end{enumerate}

\section{مسیرهایی برای تحقیقات آینده}

\subsection{استراتژی‌های انطباقی و وفقی شکل‌دهی $\rho$}

همان‌طور که در فصل بررسی چالش‌ها بحث گردید، اتکا به متغیر کنترل‌گر $\rho$ برای الگوریتم ضروری است. ما در توسعه‌های آتی در نظر داریم با برنامه‌ریزی یک ساختار هوشمند برای تنظیم نرخ فاکتور کنترل‌کننده متناسب با فرآیندهای زمانی مدل پیشروی کنیم تا قابلیت استقرار در زمان اجرای سیستم را تسهیل نماییم. برخی از جهت‌گیری‌های تخصصی شامل موارد زیر می‌باشد:
\begin{enumerate}
  \item \textbf{برنامه‌ریزی تنظیمیِ پویای پارامتر ($\rho$):} به جای یک متغیر تنظیم شده‌ی دستی، از سازوکاری جهت ارتقا یا تخفیفِ متغیر بهره خواهیم جست؛ بدین گونه که در هنگامه موفقیت‌های مکرر و تسلط ماشین (فاز استثمار) از درصدِ فاکتور $\rho$ کاسته شده و در هنگام درگیری مستمرِ ماشین و گرفتار شدن آن درون تله‌ی بهینه محلی که نرخ دریافت ارزش با فُتول روبرو می‌گردد به صورت پویا با ارتقای متغیر ضریب فشارِ فرار و اکتشاف تجربیات جدید تشدید شود. در این مسیر استفاده از مدل‌های فرا-یادگیری (\lr{Meta-Learning}) نیز در خصوص استنتاج بهینه فاکتور مفید فایده خواهد بود.
  \item \textbf{شکل‌دهی هوشمند بر پایه‌ی خطای عدم‌اطمینان:} امکان به‌کارگیری ابزارهای ارزیابی تخمین اطمینان (مانند \lr{Dropout} یا انواع مدل‌های اِنسامبل) در الگوریتم، مقدور می‌سازد تا آن دسته از تجربه‌های شکست خورده‌ای در چرخه‌ی تصادفی‌سازی اولویت یابند که سیستم درباره‌ی قطعیت خطاهای‌شان کمترین اطمینان و شناخت را در محیطِ غیرشناخته از خود نشان داده بود.
  \item \textbf{اولویت‌بندی حافظه‌ی نمونه‌گیری:} ادغام روش انتخاب هوشمند با الگوریتم‌های حافظهِ مبتنی بر اولویت (\lr{Prioritized Experience Replay})، عاملی را معرفی می‌کند تا زیرمجموعه‌های مردود شده‌ی مستعدِ جبران پس از اعمال شکل‌دهی، با فرکانس شانسِ بیشتری برای بازیابی توسط شبکه مجدداً احضار شوند.
\end{enumerate}

\subsection{گسترش به دامنه‌های پیوسته و حساس به ایمنی}

با عنایت به ماهیت پویای فضای نرم‌افزار، توسعه روش‌های حاضر بر مسائل نوین‌تر ضروری خواهد بود:
\begin{enumerate}
  \item \textbf{فضای اعمال پیوسته (\lr{Continuous Action Spaces}):} بسط معماری متدیک فعلی در بَطنِ متدهایِ قدرتمندتری نظیر سیستم‌های شیبِ سیاستی همانند متدهای \lr{PPO} و \lr{SAC} از اهداف توسعه‌ی اصلی پژوهش است تا علاوه بر شکل‌دهی پاداش تصادفی، بتوان پارامترهای بدون گسستگی (همچون تنظیم منابع سخت‌افزاری پردازش در فضای ابری) را تنظیم نمود.
  \item \textbf{اعمال محدودیت‌های تایید ایمنی (\lr{Safe Reshaping}):} برای رفع نگرانی مطرح شده در دامنه‌های حساس به ایمنی، در نظر داریم مکانیزمی‌با عنوان «شکل‌دهی ایمن» پدید آوریم که بازتولیدِ تصادفیِ نتایج شکست‌خورده تنها زمانی مُیسّر گردد که خط‌مشی پیش‌بینی‌شده در مناطق ایمن تاییدشده قرار گیرد. ترکیب با ابزارهای اعتبارسنجی صوری (همچون \lr{UPPAAL-SMC} استفاده شده در این تحقیق) از هدایت به وضعیت‌های ناایمن پیشگیری می‌کند.
  \item \textbf{گسترش به سامانه‌های چند عاملی (\lr{Multi-Agent}):} برای مقابله با فضاهایی با پویاییِ بخش‌های تطبیقی در حال تعامل با هم، گسترش متد در چارچوبی صورت می‌پذیرد که چند عامل در هماهنگی استراتژی‌های کاوش همسان به شکل‌دهی پاداش‌های منسجم برای تحقق اهداف بزرگ عمل کنند.
\end{enumerate}

\subsection{اعتبارسنجی بین‌دامنه و تحلیل‌های نظری}

جهت بررسی فراگیری این استراتژی و ایجاد پایه‌های نظری قوی‌تر اهداف زیر دنبال خواهد گردید:
\begin{enumerate}
  \item \textbf{پیکر‌بندی شبکه‌های ابری (\lr{Cloud Auto-scaling}):} به‌کارگیری این رویکرد به جهت تخصیص و توزیع مقیاس ماشین‌های مجازی، و ارزیابی در برابر کارهای مرتبط دیگر در این محیط.
  \item \textbf{هماهنگی وسایل نقلیه خودران (\lr{Autonomous Vehicles}):} آزمایش در محیط‌های شبیه‌سازی برای هماهنگ‌سازی و تطبیق الگوهای ترافیکی با محدودیت‌های ایمنی و اهداف چندگانه (زمان سفر، بهینه‌سازی انرژی و امنیت).
  \item \textbf{تضمین تئوریک فرآیند همگرایی (\lr{Theoretical Analysis}):} بررسی اینکه آیا سیگنال پاداش شکل‌دهی‌شده ویژگی انقباضی و نگاشت ثابت برای یادگیری کیو (\lr{Q-Learning}) را حفظ می‌کند، و محاسبه کران‌های خطای احتمالی انحرافِ ارزش \lr{Q} نسبت به سیاست بهینه.
  \item \textbf{کاهش حجمِ پیچیدگی نمونه‌ها (\lr{Sample Complexity Analytical Analysis}):} کمیت‌بندی تئوریک پیرامون اثربخشیِ شکل‌دهی پاداش بر تقلیل شمار نمونه‌های لازم برای دستیابی به خروجی‌های نزدیک به بهینه در قیاس با شبکه‌ی \lr{DQN} استاندارد.
\end{enumerate}

\section{نتیجه‌گیری}

\subsection{خلاصه بیان مسئله و راهکار پیشنهادی}

این رساله در قامت دستاوردی متصل بر یکی از چالش‌های بنیادین سیستم‌های خودوفقی مقاصد خود را تنظیم نموده است؛ مدیریت عدم‌قطعیت در شبکه‌هایی که درون فضاهای بزرگ و گسسته‌ی تطابق عمل می‌کنند. رویکردهای سنتی - اعم از مدل‌های مبتنی بر یادگیری ماشین، مدل‌های متکی بر جستجو و یا یادگیری تقویتی استاندارد - با محدودیت‌های شدیدی مواجه بوده‌اند؛ ویژگی‌هایی از قبیل نیاز مبرم به دیتای برچسب‌دار آفلاین، دست و پنجه نرم کردن با فقدان مقیاس‌پذیری و احتمال بالای توقف در بهینه‌های محلیِ برآمده از تغییر شرایط.

ما استراتژی \textbf{یادگیری تقویتی عمیق با شکل‌دهی تصادفی پاداش (\lr{RS-DRL})} را ارائه نمودیم؛ رویکردی نوین که مکانیسم تصادفیِ شکل‌دهی را درون جریان حافظه بازپخشِ عامل کیو (\lr{DQN}) یکپارچه می‌نماید. با بهره‌گیری از زیرمجموعه‌ای از تجربیات ناکامِ فاز گذشته ($r_i=0$) و اعمال تحول در پاداش و تبدیل آن به تجربه‌ای خوش‌بینانه ($r_i=1$)، راهکار \lr{RS-DRL} مشوق قدرتمندِ فرآیندِ اکتشاف بوده و از گرفتاری در بهینه‌های محلی جلوگیری می‌نماید، بدون آنکه کمترین الزامی‌برای تعریف و اختصاص قطعی‌ِ اهدافِ مبتنی بر دانش از پیش‌تعیین‌شدهِ دامنه، داشته باشد.

\subsection{مرور مشارکت‌ها و شواهد تجربی}

یافته‌های مستخرجِ پژوهش که ارائه‌ی پاسخی قاطع بر پرسش‌های این رساله به شمار می‌آید شامل موارد کلیدی زیر است. جدول \ref{tab:summary_contributions} خلاصه‌ای از این مشارکت‌ها و شواهد را نشان می‌دهد.

\begin{table}[h!]
\centering
\caption{خلاصه‌ای از مشارکت‌های کلیدی و شواهد تجربی}
\label{tab:summary_contributions}
\begin{tabular}{|p{4.5cm}|p{5.5cm}|p{5cm}|}
\hline
\textbf{مشارکت (\lr{Contribution})} & \textbf{شواهد (\lr{Evidence})} & \textbf{اهمیت (\lr{Significance})} \\ \hline
شکل‌دهی پاداش نوین (\lr{Novel Reward Shaping}) & توسعه الگوریتم \ref{alg:rs_drl_final}، نمایش در شکل‌های ۵ و ۶ & امکان اکتشاف بدون نیاز به دانش دامنه سیستم \\ \hline
تطبیق و وفقی‌سازی مقیاس‌پذیر & مستخرج از جداول ۳ و ۴ (نسخه‌های \lr{v1} تا \lr{v2}) & حفظ پایداریِ عملکرد با رشد گزینه‌ها از ۲۱۶ تا ۴۰۹۶ حالت \\ \hline
برتری و تفوق اثبات آماری & تحلیل جداول ۵ و ۶، معناداری ۲۵ از ۲۸ تست & دریافت \lr{$p < 0.001$} برای اکثر متریک‌های بهبود اتلاف بسته \\ \hline
\end{tabular}
\end{table}

\textbf{مشارکت اول: ابداع ماشین نوین جهت شکل‌دهی پاداش}
توسعه الگوریتم تخصیص پاداش (طراحی الگوریتم \ref{alg:rs_drl_final}) که از طریق کنترل‌کننده‌ای با ضریب مشخصِ ($\rho$) عمل نموده‌ و برخلاف رویکردهای غایت‌محور نظیر متدهای \lr{HER} هیچگونه نیازی بر تعیین هدف آستانه‌ای با درک معنایی نداشته و مستقیماً روی سیستم‌های خودوفقی قابل راه‌اندازی است.

\textbf{مشارکت دوم: یکپارچه‌سازی حلقه کنترل \lr{MAPE-K}}
طراحی زیرِسامانه‌ی متصل متشکل از چرخه ۹ مرحله‌ای که \lr{RS-DRL} را با حلقه \lr{MAPE-K} ترکیب می‌کرد؛ مجهز به ارزیابی برای تضمین اعتبارِ زمانِ اجرا با موتور کنترل \lr{UPPAAL-SMC} و مؤلفه \lr{ActivFORMS} برای اجرای وفقی.

\textbf{مشارکت سوم: دستاوردهای تجربی در مقیاس اینترنت اشیاء}
دستاورد استقرار شبکه‌ی عامل بر بستر ۳ مقیاس وسیع (فضای محدود ۲۱۶، نسخه ۱۲۹۶ و نسخه منبسط ۴۰۹۶ تصمیم ممکن). دستیابی بی‌رقیب به نرخ مجانبی $0.9905 \pm 0.0268$ (در شبکه اولیه) و $0.8983 \pm 0.0932$ (در \lr{v1.1}) که با قاطعیت ساختار کلاسیک \lr{$\epsilon$-greedy DQN} (با راندمان $0.4559 \pm 0.2233$) و مشابه‌های \lr{HER} را مغلوب ساخت. ثبت کاهش شگفت‌انگیزِ توالیِ ٪۲۲.۲۳ برای نابودی بسته‌ها و ٪۳۹.۰۷ کاهش تأخیر شبکه‌ نسبت به \lr{DLASER+} تحت اهداف \lr{TTS}. این مقادیر از طریق نتایج آزمون \lr{t-test} (موفقیت در ۸۹.۳٪ از ۲۸ مقایسه متمایز) اعتبار مطلق دریافت کرد.

\textbf{مشارکت چهارم: بینش عمیق نظری و عملی}
اثبات و آشکارسازیِ اینکه مکانیزمِ روایی برای تزریق خوش‌بینی نه تنها راه‌حلی منطقی، خروج سیستم را از گیر افتادن مصون داشته، بلکه تحلیل الگوهای متضاد درون‌سیستمی‌نظیر (تأخیر متناقض با نرخ اتلاف بسته‌ها) نشان دادیم عامل استنتاج‌هایی به وسعت دانش طراحی چند-هدفه و سازگار با رفتار سیستم هوشمند اخذ کرده است و همزمان این ظرفیت برای تغییر سایز و پدیده‌ی مقیاس‌پذیری فضای عمل بسیار کارآمد بوده است.

\subsection{سخنان پایانی}

از آنجا که سیستم‌های نرم‌افزاری به‌طور فزاینده‌ای در محیط‌هایی سرشار از عدم قطعیت، پویایی و مقیاس‌های بزرگ عمل می‌کنند، نیاز به مکانیسم‌های وفقی که بتوانند به‌طور کاملاً مستقل فرآیند یادگیری و تنظیم استراتژی را انجام دهند چشمگیرتر می‌شود. مدل هوشمند \lr{RS-DRL} یک قدم استوار در جهت این چشم‌انداز تلقی می‌شود - روشی که نه‌تنها متکی به آموزش‌های مبتنی بر موفقیت است بلکه با ظرافتی شایسته، از بطن اشتباهات پیشین درس‌های بسیار موثری اخذ می‌کند و به بیان دیگر کاستی‌های پیشین را به یک فرصت ارزنده تبدیل خواهد کرد.

اهمیت این تحقیقات مرزهای توسعه‌یافته در یک شبکه \lr{DeltaIoT} را کاملاً در هم شکسته است. هر نوع معماری که توأمان با ۳ چالش بنیادین سیستمِ نرم‌افزار (استقرار محیط پیکربندی کلان، عدم‌قطعیت‌های متغیر شرایط و نیاز تطبیق زمانِ اجرا) دست و پنجه نرم کند، با تکیه بر متد معرفی‌شده در \lr{RS-DRL} می‌تواند بر این معضل غلبه نماید، به گونه‌ای که سیستم‌های ابری، امنیتِ سایبریِ شبکه‌، کنترلِ برق هوشمند و هدایت پهپادهای یکپارچه از مصادیق این قابلیت‌ها هستند.

مأموریتِ نهایی آینده برای سیستم‌ها آن است که نه تنها مستمراً راهبریِ سیستم را هماهنگ بدانند بلکه فراتر از آن استراتژی را پیوسته غنی‌سازی نمایند – هدفی بزرگ که در شُرف تحقق پیوسته است. رساله‌ی حاضر با اتکا بر راهکار ابداعی \lr{RS-DRL} تلاش داشته خشتِ کوچکی اما ارزشمند رو به سوی این آرمان کلان مهندسی نرم‌افزار بنا کند.
\newpage
